\chapter{The Relativistic Boltzmann Equation}

\section{Special Relativity}

It is convenient to define the four-velocity of a particle as
\[
U^\alpha = \dv{x^\alpha}{\tau}
  \]

  Where $\tau$ is proper time, which is scalar invariant. The contra and co variant four velocity as

\[
(U^\alpha) = \left(\frac{c}{\sqrt{1-v^2/c^2}}, \frac{V}{\sqrt{1-v^2/c^2}} \right)
  \]

  \[
(U_\alpha) = \left(\frac{c}{\sqrt{1-v^2/c^2}}, \frac{-V}{\sqrt{1-v^2/c^2}} \right)
  \]

\subsection{Minkowski Force}

Four velocity is defined in the same manner

Then the four momentum contravariant component is given as

\[
(p^\alpha) = (\frac{E}{c}, \mathbf{p})
  \]

  Here we can derive the famous $E=mc^2$

  Taking the acceleration we get the Minkowski force's contravariant components

\[
    (K^\alpha) = (\frac{\mathbf{F} \cdot \mathbf{v}}{c(1-v^2/c^2)^{\frac{1}{2}}},\frac{\mathbf{F}}{(1-v^2/c^2)^{\frac{1}{2}}} )
 \]

\subsection{Electrodynamics in free space}

\subsubsection{Maxwell Equations}

The electromagnetic field tensor is defined through:
\[
F^{\alpha \beta} = (-\mathbf{E},c \mathbf{B})
  \]
  This tensor is antisymetric.

  With its transormation laws being
\[
    F^{'\alpha \beta} = \lambda^\alpha_\gamma \lambda_\delta^\beta F^{\gamma \delta}  \]
\[
    F^{\alpha \beta} = \bar{\lambda^\alpha}_\gamma \bar{\lambda_\delta^\beta} F^{'\gamma \delta}
  \]

  \[
\begin{gathered}
p^\mu p_\mu=-m^2c^2
\qquad\text{(mass shell)},\\
F^{\mu\nu}=\partial^\mu A^\nu-\partial^\nu A^\mu
\qquad\text{(EM field)},\\
\partial_\mu F^{\mu\nu}=\mu_0J^\nu,\qquad
\partial_\mu{}^\star F^{\mu\nu}=0
\qquad\text{(Maxwell)},\\
\frac{dp^\mu}{d\tau}=qF^{\mu\nu}u_\nu
\qquad\text{(Lorentz force)},\\
dP=\frac{d^3p}{p^0}
\qquad\text{(invariant measure)},\\
J^\mu=\sum_s q_s\int p^\mu f_s\,dP
\qquad\text{(kinetic current)},\\
p^\mu\partial_\mu f
+qF^{\mu\nu}p_\nu\frac{\partial f}{\partial p^\mu}
=C[f]
\qquad\text{(relativistic Boltzmann)},\\
T^{\mu\nu}=\int p^\mu p^\nu f\,dP
\qquad\text{(kinetic stress-energy)}.
\end{gathered}
\]

\section{Relativistic Boltzmann Equation}
\subsection{Single Non-degenerate Gas}

Note that momentum and number of particles in an volume element are scalar invariant.

We make the volume element

\[
d \mu = d^3xd^3p
  \]

  We can approximate the Jacobian to be
\[
J = 1 + \pdv{F^i}{p^2} \Delta t + O[(\Delta t)^2]
  \]

  We assume the following
  \begin{itemize}
    \item Only binary collisions are considered(reasonable if the gas is dilute)
    \item Molecular Chaos assumption
          \item One particle distribution function does not vary much over the time interval
  \end{itemize}

  Remember the momentum four vector is
\[
p^\alpha = (p^0, \mathbf{p})
  \]

  where $p^0$ is
\[
p^0 = \sqrt{\abs{p}+m^2c^2}
  \]

  where $p^0$ is the energy component and $\mathbf{p}$ is the 3-vector spacial momentum vector.


This is all pretty trivial and can be looked up later. I am going to skip ahead.

\section{Vlasov}

\subsection{The Vlasov-Maxwell system}

The Vlasov can be written in form
\[
p^\mu \pdv{f}{x^\mu} - \Gamma^\sigma_{\mu \nu}p^\mu p^\nu \pdv{f}{p^\sigma} = 0
  \]
  Here the phase space isnot a product of the two spaces but rather a sub-manifold of the tangent bundle.

  Thus the derivative is taken in the direction tangent to the manifold $g_{\mu \nu}p^\mu p^\nu = m^2 c^2$.

One can couple with Einstien's field quation by applying the stress energy tensor to be:
\[
T_{\mu \nu} = c \int p_\mu p_\nu f \sqrt{g} \frac{d^3p}{p_0}
  \]
  Creating the Vasov-Einstein system.

  \subsection{The Threshold of Black hole formation}
